%pdflatex-this
% vim:spell spelllang=en_us

\documentclass[12pt,oneside,a4paper]{report} %pdf
\usepackage[english]{babel}
\usepackage[top=25mm, bottom=15mm, includefoot, left=30mm, right=20mm]{geometry}
\usepackage[utf8]{inputenc}
\usepackage{graphicx}
\usepackage[unicode=true]{hyperref}
\usepackage{indentfirst}
\usepackage{import}
\usepackage{import}
\usepackage{moreverb}
\usepackage{multirow}
\usepackage{xcolor}
\usepackage{tabularx}
\usepackage{pdfpages}
\usepackage{ifthen}
\usepackage{cite}

\usepackage[toc,acronym,nonumberlist,style=long,sort=standard]{glossaries}
\setlength{\glsdescwidth}{0.7\linewidth}
\renewcommand*{\glsgroupskip}{}
\makeglossaries


% watermark draft version
\usepackage{draftwatermark}
\SetWatermarkText{draft \today}
\SetWatermarkLightness{0.9}


% serif font
%\renewcommand{\familydefault}{cmss}

% image numbering
\renewcommand{\thefigure}{\thesection.\arabic{figure}}
\renewcommand{\thetable}{\thesection.\arabic{table}}


%\newcommand{\Adresa}[1]{\texttt{#1}}
%\newcommand{\Cesta}[1]{\texttt{#1}}
%\newcommand{\Prikaz}[1]{\texttt{#1}}
%\newcommand{\Rozhrani}[1]{\texttt{#1}}
%\newcommand{\Nazev}[1]{\it{{#1}}}
%\newcommand{\Tabth}[1]{\textbf{{#1}}}
%\newcommand{\B}[1]{\textbf{{#1}}}
%
%\newcommand{\Th}[1]{\textbf{#1}}
%
%
\newcommand{\Uv}[1]{"#1"}

%\newcommand{\Preklad}[1]{\emph{(#1)}}

\newcommand{\Acronym}[2]{\newacronym{#1}{#1}{#2}}
%\newcommand{\Ac}[1]{\acrshort{#1}}
\newcommand{\Ac}[1]{\acrshort{#1}}

\def \BookType {Master thesis}
\def \Bookname {Measurements of virtual machines migrations}
\def \Keywords {Data center, overlay, underlay, OpenNebula, migration}
\def \Authors {Bc. Tomáš Kukrál}
\def \Supervisor {doc. Ing. Leoš Boháč, Ph.D.}
\def \Date {October 2015}
\def \Datefull {2015-11-11} % ISO8601
\def \Place {Prague}
\def \School {Czech Technical University in Prague}
\def \Faculty {Faculty of Electrical Engineering}
\def \Department {Department of Telecommunication Engineering}
\def \Programme {Communication Networks} 

\makeatletter
\def\@makechapterhead#1{%
  \vspace*{20\p@}%
  {\parindent \z@ \raggedright \normalfont
    %\ifnum \c@secnumdepth >\m@ne
    %    \huge\bfseries \@chapapp\space \thechapter
    %    \par\nobreak
    %    \vskip 20\p@
    %\fi
    \interlinepenalty\@M
    \Huge \bfseries #1\par\nobreak
    \vskip 40\p@
  }}

\makeatother

\def \Titulkadesky {
    \begin{center}
    {\LARGE \School}\\[2ex]

%    \begin{figure}[h!]
%    \begin{center}
%      {\includegraphics[width=35mm]{grafika/cvut-logo-bw-600.png}}
%    \end{center}
%    \end{figure}
    \vfill

%    {\textbf {\Huge \Bookname \\[4ex]}
        {\LARGE \bf \BookType}
%	    }

  \vspace{60mm}

    {\large \bf \Date \hfill \Authors}\\[6ex]

    \newpage
    \end{center}
    }

\def \Titulka {
    \begin{center}
    {\Huge \School}\\[2ex]
    {\Large \Faculty}\\[2ex]
    {\Large \Department}\\[2ex]

    \begin{figure}[h]
    \begin{center}
      {\includegraphics[width=4cm]{images/cvut-logo-bw-600.png}}
    \end{center}
    \end{figure}
    \vfill

    {\textbf {\Huge \Bookname \\[4ex]}
        {\LARGE \bf \BookType}
	    }

    \vfill

    {\large {\bf Author:} \hfill \Authors}\\[3ex]
    {\large {\bf Supervisor:} \hfill \Supervisor}\\[3ex]

    \newpage
    \end{center}
    }



\hypersetup{
	pdfauthor={\Authors},
	pdftitle={\Bookname},
	pdfsubject={\BookType},
	pdfkeywords={\Keywords},
	colorlinks=true,
%	hidelinks=true,
}

% environment pro abstract, poděkování a prohlášení
\newenvironment{spodnitext}[1]{
	\cleardoublepage
	\null
	\vfill
	\section*{#1}
	}{
	\vspace{10mm}
	}


\begin{document}
% acronyms
\subimport{./}{acronyms}


\catcode`\-=12

%\begin{titlepage}
%	\Titulkadesky
%\end{titlepage}

% vypnout watermark

% prázdná stránka
%\newpage
%\null
%\clearpage

\begin{titlepage}
%	\Titulka
\end{titlepage}
\SetWatermarkText{}

\pagenumbering{Roman}

%prohlášení
\begin{spodnitext}{Proclamation}
I declare that this thesis was mady by myself with assistance of my supervisor. All parts taken over word by word from literature or other publications are referenced and identified. I approve publishing this thesis or any part of it with referencing author of original text.

\vspace{10mm}
In~\Place~at~\Datefull \hfill \hbox to 65mm{\tiny\dotfill}
\end{spodnitext}

% acknowlegement
%\begin{spodnitext}{Acknowledgement}
%Tímto bych chtěl poděkovat mému vedoucímu Ing. Pavlovi Bezpalcovi, Ph.D. za pomoc s~realizací práce a za cenné rady k~tématu. Bez jeho pomoci bych nemohl vytvořit tuto práci. Vážím si toho, že mi umožnil pracovat na tomto tématu.
%\end{spodnitext}


%zadání
%\includepdf[pages=-]{grafika/zadani.pdf}

% abstrakt
\begin{spodnitext}{Abstract}
%Práce se zabývá použitím protokolu \Zkratka{IPv6}{Internet Protocol version 6} na sítových prvcích Juniper s~operačním systémem \Zkratka{JUNOS}{Juniper Operating System}, zejména řadě SRX. Zaměřuje se na praktické problémy s~nasazením v~síti a problematikou přidělování \Zkratka{IP}{Internet Protocol} adres a síťových parametrů.\\
%\textbf{Klíčová slova:} \Zkr{IPv6}, \Zkr{JUNOS}, Juniper, \Zkratka{DHCPv6}{Dynamic Host Configuration Protocol version 6}, \Zkratka{SLAAC}{Stateless Address Autoconfiguration}
%
%\vspace{1cm}
	This master thesis deals with testing data networks in cloud environment. Techniques of inter and intra-cloud networks are described in theoretical part as well as virtual machine migrations.
	Practical part brings methodology and framework for testing virtual machine migration. Measurements are performed at OpenNebula cloud environment with {KVM} virtual machines.
\end{spodnitext}


%obsah
\tableofcontents
\cleardoublepage

\pagestyle{plain}
\pagenumbering{arabic}

\chapter{Introduction}
\subimport{./}{introduction}

\chapter{Theoretical part}
	\section{Virtualization}
		\subsection{Types of virtualization}
		\subsection{Levels of virtualization}
			\subimport{./theoretical/}{theoretical/levels}
		\subsection{Advantages of virtualization}
	\section{Cloud computing}
		\subimport{./theoretical/}{theoretical/introduction}

		\subsection{Network in cloud}
			% overlays
		\subsection{Storage in cloud}
			% share-nothing, centralized storage

		\subsection{Orchestration software}
			\subsubsection{OpenNebula}

		%\section{Migration of VMs}

	\section{Distributed datacenter}

	\pagebreak

	%Sharing resources - maximization of effectiveness
	%Public/private/hybrid
	%metrics
	%SOA - service oriented architecture
	%Routing requests to virtual machines
	%Network environment
	%	overlays
	%Data storage
	%Migration
	%	cold/live
	%Measurement methods
	%Economy
	%	capex -> opex
	%	pay as you goo
\chapter{Practical part}

	Methodology overview

	Framework

	Results

%zdroje
%	\subimport{/}{zdroje}


%seznam obrázků
\newpage
\phantomsection \label{listoffig}
\addcontentsline{toc}{chapter}{List of Figures}
\listoffigures

%seznam tabulek
\newpage
\phantomsection \label{listoftab}
\addcontentsline{toc}{chapter}{List of Tables}
\listoftables

%seznam pojmů
%\printglossary[title=Acronyms,toctitle=Acronyms used]

%seznam zkratek
\printglossary[type=acronym,title=List of Abbreviations,toctitle=List of Abbreviations]

% přílohy
%\appendix
%\chapter{Příloha}

%\section{Konfigurace prvků}
%\label{priloha_konfigurace}
%{\scriptsize 
%\subsection{R1}
%	\verbatimtabinput{prakticka/conf/R1_P2.conf}
%\subsection{R2}
%	\verbatimtabinput{prakticka/conf/R2_P2.conf}
%\subsection{R3}
%	\verbatimtabinput{prakticka/conf/R3_P2.conf}
%}


%\cleardoublepage
%\section{Úvod do JUNOSu}
%\label{priloha_manual}
%V~následující části je vložen manuál, který popisuje základní ovládání a konfiguraci operačního systému \Zkr{JUNOS}. Slouží k~doplnění praktické části, ve které jsem se soustředil především na samotnou konfiguraci, zatímco základní nastavení zařízení a ovládání je popsáno v~této příloze.
%\includepdf[pages=-]{../branches/junos_manual/manual_import.pdf}

%\cleardoublepage
%\section{Address autoconfiguration in IPv6 networks}
%\includepdf[pages=-]{../branches/clanek/tsp12/clanek_import.pdf}


\end{document}
